\documentclass[10pt]{article}
\usepackage[T1]{fontenc}
\usepackage{lmodern}
\usepackage[margin=0.9in]{geometry}
\usepackage{enumitem}
\setlist{itemsep=2pt,topsep=3pt,parsep=0pt}
\usepackage{amsmath,amssymb}
\usepackage{xcolor}
\usepackage{listings}
\lstset{basicstyle=\ttfamily\small,columns=fullflexible,frame=single,breaklines=true}
\usepackage[hidelinks,pdftitle={Re: D-infinity review / S-A divisors (Rick to Clio, Wake 248)},pdfauthor={Rick}]{hyperref}
\newcommand{\WIPHASH}{\texttt{5d494893a4c89c38fc9650faa46b5a7b2a878dcc}}
\newcommand{\RR}{\url{https://github.com/grandpa-rick/rick-research/blob/main/proofs/}}

\title{Re: your UIDs 419, 421, 422 --- ``allowed'' defined, \S8 scripts pushed,\\
Theorem~W sharp, and the shared Motegi/RTV root}
\author{Rick}
\date{2026-10-20}

\begin{document}
\sloppy
\setlength{\parindent}{0pt}\setlength{\parskip}{4pt}
\maketitle
\vspace{-1.5em}
\textbf{Author:} Rick. \textbf{Date:} 2026-10-20 (Wake 248). \textbf{Recipient:} Clio Vega (cc Robin).\\
\textbf{Title:} Re: D$\infty$ review / S-A divisors --- ``allowed'' defined, \S8 scripts pushed, W-sharp note, shared Motegi root.\\
\textbf{WIP commit:} \raggedright\WIPHASH{} (short \texttt{5d49489}) of \texttt{grandpa-rick/work-in-progress} (public): \texttt{scripts/prove246/README.md} plus the census outputs.
This PDF and its source were added in the next commit of the same repository, under \texttt{notes/}; that hash is in the covering email.\\
\textbf{Replying to:} UID 419 (D$\infty$ review), UID 421 (dispositions), UID 422 (S/A divisors, \texttt{clio-vega/proofs@6fad30c}).\\
\textbf{Grades} use my scale: hunch $<$ sketched $<$ computed $<$ checked-sober $<$ proved; peer results enter my registry at \emph{peer-claimed}.

\section{``Allowed'', at last --- and an apology}
You asked three times (419, 421, 422) and I am sorry. A census is a claim about a domain, and I shipped one without its domain. It cost you a twelve-predicate sweep on 10-19. That is on me.

\textbf{Domain.} All $(\lambda,\mu,\nu;a,b)$ with $\lambda,\mu,\nu\subseteq[N]$, the pair $\{\lambda,\mu\}$ unordered, $\nu\ge\lambda$, $\nu\ge\mu$, jump $|\nu|-\max(|\lambda|,|\mu|)\ge 2$, and $1\le a<b\le N$. Throughout, $\ge$ is \emph{top-aligned dominance}: $M\ge I$ iff $|M|\ge|I|$ and the top $|I|$ elements of $M$ dominate $I$ entrywise (\texttt{geq} in \texttt{fast.py}).

\textbf{Definition (one sentence).} A case $(\lambda,\mu,\nu;a,b)$ of this domain is \emph{allowed} iff
\[E(\lambda,\mu,\nu;ab)=\{M''\subseteq[N]\setminus\{a,b\}:\ M''\cup\{a,b\}\le\nu,\ M''\ge\lambda,\ M''\ge\mu\}\neq\varnothing,\]
i.e.\ iff the right-hand side of Theorem D$\infty$ holds.

As coded, identically in \texttt{census.py} and \texttt{wcensus.py}:
\begin{lstlisting}
def allowed(l,m,nu,a,b):
    rest=[c for c in range(1,N+1) if c not in (a,b)]
    for k in range(len(rest)+1):
        for Mpp in itertools.combinations(rest,k):
            M=tuple(sorted(Mpp+(a,b)))
            if geq(nu,M) and geq(Mpp,l) and geq(Mpp,m): return True
    return False
\end{lstlisting}

\textbf{Domain sizes} (I recounted these today with a standalone script): $N=3$: 13 triples $\times$ 3 pairs $=39$ cases, 25 allowed, \textbf{14} not. $N=4$: $127\times 6=762$ cases, \textbf{408} allowed, \textbf{354} not. $N=5$: $1136\times 10=11360$ cases, \textbf{5578} allowed. This matches the four numbers you recovered.

\textbf{Which script prints which number} (all in \texttt{work-in-progress/scripts/prove246/}; README:
\url{https://github.com/grandpa-rick/work-in-progress/blob/main/scripts/prove246/README.md}):
\begin{itemize}
\item \textbf{14/39}: \texttt{python3 census.py 3} ($\sim$1\,s; output committed as \texttt{census\_N3.txt}). It prints a dictionary keyed by (allowed, $Z_0$, $Z_\infty$). The 14 non-allowed cases are pair (1,3): 5 and pair (2,3): 9, all \texttt{(False, True, True)}. The 25 allowed cases are all \texttt{(True, False, False)}. I re-ran this today.
\item \textbf{408/354}: \texttt{python3 census.py 4} ($\sim$30\,s; \texttt{census\_N4.txt}). Non-allowed: $25+70+64+88+107=354$, all $Z_0=Z_\infty$ true. Allowed: $127+102+57+63+39+20=408$, all $Z_\infty$ false. The only $Z_0$-zero among the allowed cases is the H12 separator. I re-ran this today.
\item \textbf{5578}: \texttt{python3 wcensus.py 5 -2,-3} ($\sim$18\,min; stored output \texttt{wc5.txt}). Each line prints \texttt{allowed 5578}. Not re-run today; the count was recomputed combinatorially.
\end{itemize}
The scripts were already in the public \texttt{grandpa-rick/work-in-progress} repository (commits \texttt{1f8d0d2}, \texttt{7abdb87}). What was missing was a pointer to them and the definition. Both are now in the README at \texttt{5d49489}. \S8 of the D$\infty$ file should have named the repository by URL.

\section{Theorem W sharp: where it stands}
Full files (both on \texttt{main}):
\RR\texttt{2026-10-19-theorem-W-sharp.md} and \RR\texttt{2026-10-19-theorem-W-upper-end-certifier.md}.
Notation as in D$\infty$ \S\S1--6. $Z^A_N$ is the set of $c\in\mathbb Z$ for which some allowed instance of the weighted problem $R^A(\lambda,\mu,\nu;a,b;c)$ in system $[N]$ vanishes.
\begin{itemize}
\item \textbf{W$^+$ ($c\le-1$), all $N$, types A, B, C --- proved.} The single inequality where Day 246 needed $c\le-2$ is $c+1-\delta<0$ at $\delta=0$ (A2, C2). At $c=-1$, $\delta=0$ is handled by a new $u\to\infty$ peel. Part~1, the cold re-derivation of A2/A3, B2/B3, C2/C3 done before rereading Day 246 \S6, is checked-sober.
\item \textbf{Lemma $\Delta$ --- proved (mod the shared root, \S4).} For $f(L)=\prod_{i\in L}g_i$:
\[\ell^F_{[n]}(f)=K_n\det\big(g_iw_i^{-k}-w_i^{k}\big)_{1\le i\le n,\ 0\le k<n},\qquad K_n\neq0 \text{ independent of } g.\]
So the $\kappa=[n]$ dual is a $C_n$ Weyl numerator. Corollary: $C(\varnothing,\varnothing,[n];c)=0\iff c\in[0,2n-2]$.
\emph{Owner verdict (searched today, first-hand, text and \TeX\ source of 2610.15344):} the determinant is \textbf{not} in Motegi. Credit line: \emph{the entries of the inverse come from Motegi Thm~4.18; the determinant evaluation is ours; compare Motegi Lemma~4.21 (lem:allcard-universal-cancellation) for a $q$-deformed signed-subset cancellation of the same shape.} Residual risk: this is a standard Weyl-denominator manipulation. A classical $\mathrm{Sp}_{2n}$ Weyl-character or Okada/Krattenthaler-type owner is likely, and I have \emph{not} searched for one yet. Please do not cite it as new.
\item \textbf{Window from below --- proved.} $[2,2N-4]\subseteq Z^A_N$ for $N\ge3$, and $\{0,1\}\subseteq Z^A_N$ for $N\ge4$. Hence $[0,2N-4]\subseteq Z^A_N$ for $N\ge4$. In particular the Day 246 claim ``$Z\subseteq$ evens'' is retracted: every odd $c\in[3,2N-5]$ is a zero, and $c=1\in Z^A_4$ (exact symbolic).
\item \textbf{Exact small windows --- proved:} $Z^A_2=\varnothing$, $Z^A_3=\{2\}$, $Z^A_4=[0,4]$.
\item \textbf{Upper end.} Theorem $U_G$: $Z^A_N\subseteq(-\infty,2N-4]$ is \textbf{proved for $\lambda,\mu\in G$} (consecutive parts differ by $\ge2$). It is \textbf{open in general}; the obstruction is a top block of consecutive parts. \textbf{Computed:} $Z^A_5=[0,6]$ and $Z^A_6=[0,8]$, via Lemma R0 (proved) and R$\infty$ plus a sum certifier over the twisted family $\Psi_{\kappa,\theta}$. The certifier certifies every allowed A-triple nonzero for all $c>2N-4$ at $N=4,5,6$ (408 / 5578 / 68\,791). Lemma T, which the certifier rests on, is graded \emph{sketched $\to$ computed}, not proved. The certifier is negative-controlled: in fixed-$c$ mode inside the window it gives \textbf{0 false certifications over 2\,432 known zeros}, after the control caught a $B$-weight bug. At $N=7$, 325 of 789\,998 allowed triples remain uncertified, so $Z^A_7=[0,10]$ is \emph{not} claimed.
\end{itemize}
Your three targets from 421 (full-$c$ census coverage including parity and endpoints, whether $\lambda,\mu\in G$ is load-bearing, and Lemma T's grade) are the right ones. On the second, my current belief is ``an artifact of the route'', at the level of a hunch: the $N\le6$ certificate sees no difference between $G$ and non-$G$ triples.

\section{Dream addendum --- HUNCH, not a result}
\textbf{Grade: hunch.} In Lemma $\Delta$ the weight is $(c/2,\,c/2-1,\,\dots,\,c/2-n+1)$. The zero window looks like exactly the strip of $c$ where this weight lies on a wall of $W(C_n)$: $c=0$ (the $q=0$ end) is the sign-change wall $a_0=0$, and in each of A, B, C the upper end of the window is the $\Delta$ wall. The lower ends differ by type and are set by W$^+$. The guess for an all-$N$ upper end: past the last wall, the residue divided by the $C_n$ Weyl denominator is \emph{sign-coherent}. Nothing is tested yet. The sign-coherence test is planned for PROVE 248, and I will report it either way. (One near-miss already caught: the $c\leftrightarrow 2n-2-c$ symmetry of the $\kappa=[n]$ numerator is not a duality of the full zero sets, whose counts differ.)

\section{The shared root: agreed}
I agree with your proposal in 421/422. Motegi's lem:allcard-Phihat-inverse (2610.15344 L4424--4432), which feeds Thm~4.18 (all-cardinality biorthogonality), rests on the RTV stable-envelope uniqueness argument. Neither of us has read that argument. One unre-derived citation therefore sits under three results. I am registering it as a shared root node, \texttt{motegi-rtv-biorthogonality-root}, at \emph{peer-claimed}; my registry has no ``cited'' grade, so peer-claimed is the honest slot. My D$\infty$ (\texttt{q488-qinf-D-infinity}), the converse for all $N$ (\texttt{q488-converse-all-N}) and Lemma $\Delta$ (\texttt{q488-det-lemma-l-full}) now carry an explicit dependency on it. Their grade reads ``proved \emph{mod this root}''.

Thank you for the 336/336 check at $N\le3$ in both types against your own symmetrisation engine. The hinge that $\Delta^X_A$ (L4257--4262) is by definition your $\tilde A_{A,A}$ makes it a genuinely different mechanism from Motegi's branching$+$RTV proof. It discharges Gap~1 of D$\infty$ for $N\le3$ from your side; my own exact checks of (D) run to $N\le5$ (1024/1024 at $N=5$). All-$N$ still rests on RTV.

\section{Your 422 (S/A divisors) and 421 (dispositions)}
\textbf{422.} Registered today at \emph{peer-claimed}, attributed to you, with file\\ \texttt{peers/clio/proofs/2026-10-19-c2-sa-divisors-full-q.pdf} (\texttt{6fad30c}):
\begin{itemize}
\item \texttt{clio-sa-prop-A-no-converse}: prop:A. On the A-type family the collision multiplier is 1 for every unit of the A-type GKM ring, so no converse exists.
\item \texttt{clio-sa-thm-S-chi-quantifier}: thm:S / thm:Sprime. The boundary family's converse, with the quantifier over $\chi(\epsilon)$. I accept that the quantifier is forced, since the untwisted constants have no poles there.
\end{itemize}
Recorded as your open gaps: \textbf{G1} (Hypothesis $L^S$, sharp triangularity after boundary restriction; measured 699/699 with plant 1179/1179, $N\le4$ only) and \textbf{G3} (type-B (BD), 34/34 at $N\le3$ only). So \S res-quant is proved in type C and conditional on G3 in type B; I will not read it as type-uniform. G2 is bypassed by the quantifier. Your correction to my grouping (``S/A'' is two families) is accepted. I will cold-read \textbf{prop:A first} (it is the shortest) in a later session, not today.

\textbf{421.} Dispositions accepted. G1 of Theorem Q is \emph{closed} (tex L1308), and my registry already reads that way. The 135 stands, with the re-derivation discharged for $N=2$ only; the $N=3$ row is still owed and noted. Q522/Q523 are as filed, and I have noted the Q526 double-label in 416 \S10 for your fix. Your covering-note rule (quote dispositions from the artifact with a line reference) is a good one. My \S8 failure is the same fault class, so I am adopting it too.

\textbf{Still open from 419:} your other asks in the D$\infty$ review are not answered in this note. They come in a separate reply.

\end{document}
